%------------------------------------------------------------------------------%
\begin{question}
  Considerando a função
  \(
    f(x, y) = x^3 + 3xy + y^3
  \)

  \smallskip
  a) Calcule o gradiente de $f$

  \smallskip
  b) Calcule a hessiana de $f$

  \smallskip
  c) Encontre todos os pontos críticos de $f$

  \smallskip
  d) Classifique cada ponto crítico de $f$
\end{question}

%------------------------------------------------------------------------------%
\begin{resolution}{a)\; Gradiente de $f$}
  Derivadas parciais de $f$

  \[
    \pause
    \D{f}{x}
    \pause
    \;=\; \D{}{x}\left[x^3 + 3xy + y^3\right]
    \pause
    \;=\; 3x^2 + 3y
  \]

  \[
    \pause
    \D{f}{y}
    \pause
    \;=\; \D{}{y}\left[x^3 + 3xy + y^3\right]
    \pause
    \;=\; 3x + 3y^2
  \]

  \pause
  então
  \[
    \pause
    \nabla f = \vetor{3x^2 + 3y}{3x + 3y^2}
  \]
\end{resolution}

%------------------------------------------------------------------------------%
\begin{resolution}{b)\; Hessiana de $f$}
  Derivadas parciais de segunda ordem de $f$
  \[
    \pause
    \DS{f}{x}
    \pause
    \;=\; \D{}{x}\left[3x^2 + 3y\right]
    \pause
    \;=\; 6x
  \]
  \[
    \pause
    \D{f}{y}
    \pause
    \;=\; \D{}{y}\left[3x + 3y^2\right]
    \pause
    \;=\; 6y
  \]
  \[
    \pause
    \DM{f}{x}{y}
    \pause
    \;=\; \D{}{x}\D{f}{y}
    \pause
    \;=\; \D{}{x}\left[3x + 3y^2\right]
    \pause
    \;=\; 3
  \]

  \pause
  então
  \[
    \pause
    H = \left(\begin{array}{cc}
      6x & 3 \\
      3  & 6y
    \end{array}\right)
  \]
\end{resolution}

%------------------------------------------------------------------------------%
\begin{resolution}{c)\; Pontos Críticos de $f$}
  A função é um polinômio, então possui derivadas em todos os pontos do plano

  \pause
  Os pontos críticos são apenas onde as derivadas parciais são zero,
  \(\nabla f = 0\)

  \[
    \pause
    3x^2 + 3y = 0
    \qquad\text{e}\qquad
    \pause
    3x + 3y^2 = 0
  \]
  \pause
  simplificando
  \[
    \pause
    x^2 + y = 0
    \qquad\text{e}\qquad
    \pause
    x + y^2 = 0
  \]
\end{resolution}

%------------------------------------------------------------------------------%
\begin{resolution}{c)\; Pontos Críticos de $f$}
  \onslide<+->{Isolando $y$ na primeira equação}
  \begin{align*}
    \onslide<+->{x^2 + y & = 0} \\
    \onslide<+->{y & = -x^2}
  \end{align*}
  \onslide<+->{e substituindo na segunda}
  \begin{align*}
    \onslide<+->{x + y^2                 & = 0 \\}
    \onslide<+->{x + \left(-x^2\right)^2 & = 0 \\}
    \onslide<+->{x + x^4                 & = 0 \\}
    \onslide<+->{x(1+ x^3)               & = 0}
  \end{align*}
  \onslide<+->{As soluções dessa equação são $x=0$ ou $x-1$}
\end{resolution}

%------------------------------------------------------------------------------%
\begin{resolution}{c)\; Pontos Críticos de $f$}
  Como \(y = -x^2\)

  \pause
  \medskip
  se $x=0$, então $y=0$

  \pause
  \medskip
  se $x=-1$, então $y=-1$

  \pause
  \bigskip
  os pontos críticos são
  \[
    (x_1, y_1) = (0, 0)
    \pause
    \qquad\text{e}\qquad
    (x_2, y_2) = (-1, -1)
  \]
\end{resolution}

%------------------------------------------------------------------------------%
\begin{resolution}{d)\; Classificando os Pontos}
  Considerando o ponto \((x_1, y_1) = (0, 0)\)
  \[
    \pause
    f_{xx}(0, 0)
    \pause
    = \left(6x\right)\evalat{(0, 0)}{}
    \pause
    = 0
  \]
  \[
    \pause
    f_{yy}(0, 0)
    \pause
    = \left(6y\right)\evalat{(0, 0)}{}
    \pause
    = 0
  \]
  \[
    \pause
    f_{xy}(0, 0) = 3
  \]
  \[
    \pause
    D_1
    \pause
    \;=\; f_{xx}(0, 0)f_{yy}(0, 0) - f^2_{xy}(0, 0)
    \pause
    \;=\; 0 \times 0 - 3^2
    \pause
    \;=\; -9
    \pause
    \;<\; 0
  \]

  \pause
  Portanto, o ponto $(0, 0)$ é um ponto de sela.
\end{resolution}

%------------------------------------------------------------------------------%
\begin{resolution}{d)\; Classificando os Pontos}
  Considerando o ponto \((x_2, y_2) = (-1, -1)\)
  \[
    \pause
    f_{xx}(-1, -1)
    \pause
    = \left(6x\right)\evalat{(-1, -1)}{}
    \pause
    = -6
  \]
  \[
    \pause
    f_{yy}(-1, -1)
    \pause
    = \left(6y\right)\evalat{(-1, -1)}{}
    \pause
    = -6
  \]
  \[
    \pause
    f_{xy}(-1, -1) = 3
  \]
  \[
    \pause
    D_2
    \pause
    \;=\; f_{xx}(-1, -1)f_{yy}(-1, -1) - f^2_{xy}(-1, -1)
    \pause
    \;=\; (-6)(-6) - 3^2
    \pause
    \;=\; 36 - 9
    \pause
    \;=\; 25 > 0
  \]

  \pause
  Portanto, o ponto $(-1, -1)$ é um máximo ou mínimo local

  \pause
  Como \(f_{xx}(-1, -1) = -6 < 0\)

  \pause
  o ponto é um ponto de máximo local.
\end{resolution}

%------------------------------------------------------------------------------%
